\caption{Homogeneous Sub-Sample Analysis}
\label{tab:homogeneous_subsample}
\begin{adjustbox}{max width=\textwidth}
\begin{threeparttable}
\small
\begin{tabular}{lc}
\toprule
Sub-sample & Result \\
\midrule
\multicolumn{2}{c}{\textit{Analysis infeasible}} \\
\bottomrule
\end{tabular}
\begin{tablenotes}
\small
\item \textit{Notes:} The BEC item classification uses numeric codes without
standardised product-level descriptions (e.g., ``fuel,'' ``A4 paper,''
``generic medicines''). Because item codes are hierarchical and do not map
one-to-one to homogeneous commodity groups, we were unable to reliably identify
purely homogeneous sub-samples. The item fixed effects in the baseline
specification absorb time-invariant item heterogeneity and serve a similar
purpose. Future work using more granular product catalogues may enable this
analysis.
\end{tablenotes}
\end{threeparttable}
\end{adjustbox}
